% Generated draft: requires mathematical and source review.
\documentclass[11pt]{article}
\usepackage[margin=25mm]{geometry}
\usepackage{fontspec}
\setmainfont{Latin Modern Roman}
\usepackage{amsmath,amssymb,mathtools,unicode-math}
\setmathfont{Latin Modern Math}
\usepackage{enumitem,needspace,xurl,hyperref}
\setlength{\parindent}{0pt}
\setlength{\parskip}{5pt}
\setlength{\emergencystretch}{3em}
\newcommand{\N}{\mathbb{N}}
\newcommand{\Z}{\mathbb{Z}}
\newcommand{\Q}{\mathbb{Q}}
\newcommand{\R}{\mathbb{R}}
\newcommand{\C}{\mathbb{C}}
\newcommand{\F}{\mathbb{F}}
\newcommand{\A}{\mathbb{A}}
\newcommand{\PP}{\mathbb{P}}
\newcommand{\E}{\mathbb{E}}
\newcommand{\eps}{\varepsilon}
\newcommand{\dd}{\,\mathrm{d}}
\begin{document}
\section*{UnsolvedMath ID: 30001921}
\textbf{Source problem code:} catalogue entry
\subsection{Definitions and mathematical statement}
Let \(S\) and \((S^\vee,F_{S^\vee})\) be a mirror pair of the kind intended by the supplied source. Here \(S^\vee\) denotes the mirror space and \(F_{S^\vee}\) its named additional datum; the record does not define this datum, the mirror construction, or the hypotheses on these spaces. Write \(HH^k(S)\) for the invariant called \emph{Hochschild homology} in the source, retaining its superscript notation, and \(H^k(S^\vee,\C)\) for the displayed cohomology with complex coefficients.

The source conjecture asks whether, for every such mirror pair and each intended degree \(k\),
\[
HH^k(S)\cong H^k(S^\vee,\C).
\]
The degree range is not specified in the record. Neither the category defining Hochschild homology, its grading convention, nor the cohomology theory and support convention on the right are specified. In particular, the superscript in \(HH^k\) must not silently be interpreted as Hochschild cohomology or replaced by a homological subscript. The role of \(F_{S^\vee}\), which does not occur explicitly on the right-hand side, also remains unspecified. Thus this display preserves the supplied conjectural comparison but is not a fully determined mathematical assertion until these qualifications are resolved. No naturality or compatibility with additional structures is stated.
\subsection{Short English statement}
The conjecture compares the source's degree-\(k\) Hochschild homology of \(S\) with degree-\(k\) complex cohomology of its mirror space \(S^\vee\). The mirror construction, hypotheses, grading, and cohomology conventions require clarification. Supplied dated curation labels the problem partially solved and reports results for curves and one quintic-surface example; those claims are not independently verified here.
\subsection{Sources}
\begin{itemize}
\item Supplied local UnsolvedMath registry record: permanent ID \(30001921\), exact source problem code \texttt{catalogue entry}, title \emph{Hochschild Homology of Landau--Ginzburg Mirrors}, in the Oberwolfach Reports Open Problems collection. The proposed year is recorded as 2011; no proposer is named.
\item Supplied source citation: \emph{Complex Algebraic Geometry (2012)}, Oberwolfach Reports, pp.~2669--2733, DOI \texttt{10.4171/owr/2011/47}. Original OWR report: \url{https://doi.org/10.4171/owr/2011/47}. The background describes Oberwolfach Reports as documenting mathematical workshops and their open-problem discussions.
\item Supplied literature reference, labeled Gross--Katzarkov--Ruddat: \url{https://arxiv.org/abs/1202.4042}.
\item Literature assessment checked at \texttt{2026-08-21}: this is supplied dated curation, not independently verified here. It labels the problem \texttt{partially\_solved}, reports that the conjecture is proved for curves and an example of a quintic surface in the cited work, and says that no general theorem covering the stated mirror construction was found. The record explicitly requests independent verification of time-sensitive status and cited claims.
\item The registry labels the statement \texttt{corrected\_verified} and reports that its mathematical content was checked against the official Oberwolfach report, with notation, context, and question form normalized without changing the source-level claim. This is a supplied verification assertion, not an independent source check performed for this draft.
\end{itemize}
\end{document}
